\begin{center}\begin{tikzpicture}[font=\small,>=stealth,line join=round]
\draw[gray!40,rounded corners=3pt] (0,0) rectangle (14.2,8.3);
\node[text width=12.6cm,align=center] at (7.1,7.65) {偶函数 + 轴 $x=1$};

\begin{scope}[shift={(3.2,4.3)},x=1.35cm,y=1.15cm]
\draw[->,gray] (-1.7,0)--(5.25,0) node[right]{$x$};
\draw[->,gray] (0,-1.45)--(0,1.65) node[above]{$y$};
\draw[blue,thick,domain=-1.5:5,samples=240] plot(\x,{cos(deg(pi*\x))});
\draw[violet,dashed] (1,-1.4)--(1,1.45) node[above]{$x=1$};
\fill[orange!85!black] (-0.25,{cos(deg(pi*-0.25))}) circle(2.2pt);
\fill[orange!85!black] (0.25,{cos(deg(pi*0.25))}) circle(2.2pt);
\fill[orange!85!black] (2.25,{cos(deg(pi*2.25))}) circle(2.2pt);
\fill[orange!85!black] (4.25,{cos(deg(pi*4.25))}) circle(2.2pt);
\foreach \v in {-1,1,2,3,4,5}{\node[below=4pt] at (\v,-1.5){$\v$};}\end{scope}
\node at (7.1,6.85) {$f(x)=\cos(\pi x),\quad a=1$};
\node[align=center] at (7.1,1.3) {$f(x+2)=f(x)$；$2$ 是一个正周期。};
\node[text width=12.6cm,align=center] at (7.1,.65) {橙点的横坐标为 $-\frac14,\frac14,\frac94,\frac{17}4$；曲线仅是满足全部条件的一个例子。};
\end{tikzpicture}\end{center}
